\documentclass[11pt, a4paper]{article}

% Geometria e margini della pagina
\usepackage[a4paper, margin=2.5cm]{geometry}

% Pacchetti standard per fisica, formule ed elenchi
\usepackage{amsmath, amssymb}
\usepackage{enumitem}
\setlist[itemize]{label=-}
\usepackage{booktabs}

% Collegamenti ipertestuali
\usepackage[colorlinks=true, linkcolor=blue, urlcolor=blue]{hyperref}

% Dati per l'intestazione
\title{\textbf{Verifica della resistenza interna di una lampadina ad incandescenza}}
\author{
    \textbf{Studenti: Bressan Eddi, Feoli Francesco, Rosa Carlotta}
}
\date{\today}

\begin{document}

\maketitle

\section{Obiettivo dell'Esperienza}
Descrivere brevemente lo scopo della prova (ad esempio la verifica sperimentale di una legge fisica o la misura di una costante, es. $g$).

\section{Cenni Teorici}
Esposizione sintetica delle equazioni e dei principi fisici alla base dell'esperienza:
\begin{equation}
    T = 2\pi \sqrt{\frac{l}{g}}
\end{equation}
dove $T$ rappresenta il periodo di oscillazione, $l$ la lunghezza del filo e $g$ l'accelerazione di gravità.

\section{Strumentazione Utilizzata}
\begin{itemize}
    \item \textbf{Cronometro digitale:} sensibilità $\pm 0{,}01\,\text{s}$
    \item \textbf{Metro a nastro:} sensibilità $\pm 1\,\text{mm}$
    \item \textbf{Massa e filo inestensibile}
\end{itemize}

\section{Dati Sperimentali e Misure}
I valori rilevati durante le prove sono riportati nella seguente tabella riassuntiva:

\begin{table}[h!]
    \centering
    \begin{tabular}{cccc}
        \toprule
        \textbf{Misura n.} & \textbf{Lunghezza $l$ (m)} & \textbf{Tempo $10T$ (s)} & \textbf{Periodo $T$ (s)} \\
        \midrule
        1 & $0{,}800 \pm 0{,}001$ & $17{,}95 \pm 0{,}01$ & $1{,}795 \pm 0{,}001$ \\
        2 & $0{,}800 \pm 0{,}001$ & $17{,}92 \pm 0{,}01$ & $1{,}792 \pm 0{,}001$ \\
        3 & $0{,}800 \pm 0{,}001$ & $17{,}98 \pm 0{,}01$ & $1{,}798 \pm 0{,}001$ \\
        \bottomrule
    \end{tabular}
    \caption{Raccolta delle misure sperimentali.}
\end{table}

\section{Elaborazione dei Dati e Propagazione degli Errori}
Calcolo del valore medio e stima dell'incertezza associata:
\begin{equation}
    \bar{T} = \frac{1}{N} \sum_{i=1}^{N} T_i
\end{equation}
Ricavando l'accelerazione di gravità dalla formula inversa:
\begin{equation}
    g = \frac{4\pi^2 l}{T^2} = (9{,}81 \pm 0{,}05)\,\text{m/s}^2
\end{equation}

\section{Conclusioni e Discussione dei Risultati}
Breve commento sulla compatibilità del risultato ottenuto rispetto al valore atteso teorico, analizzando eventuali fonti di errore sistematico o casuale.

\end{document}